\documentclass[11pt]{article}
\usepackage[margin=1in]{geometry}
\usepackage{amsmath,amssymb,amsthm,mathtools,booktabs}
\usepackage[T1]{fontenc}
\usepackage[colorlinks=true,linkcolor=blue,urlcolor=blue,citecolor=blue]{hyperref}
\newtheorem{theorem}{Theorem}
\newtheorem{proposition}[theorem]{Proposition}
\newcommand{\E}{\mathbb{E}}
\newcommand{\R}{\mathbb{R}}
\newcommand{\Var}{\operatorname{Var}}
\newcommand{\tr}{\operatorname{tr}}
\newcommand{\norm}[1]{\left\lVert #1\right\rVert}
\title{H-JEPA versus Flat JEPA\\Conditional Guarantees, Gains, and Losses}
\author{Working research note}
\date{October 20, 2026}
\begin{document}
\maketitle
\begin{abstract}
We formulate a mathematical comparison around arXiv:2610.06805v1. The central tradeoff is between information discarded by abstraction and the learning or planning difficulty that abstraction removes. We prove an exact prediction-risk decomposition, a finite-sample gain/loss threshold, a linear predictive-feature selection result, and a conditional planning-regret bound. These are standard mathematical arguments and explicitly simplified models, organized as a research foundation. They do not establish convergence or universal superiority of the complete trained neural architectures. Novelty remains unassessed.
\end{abstract}

\paragraph{Formal verification status.}
The companion Lean 4 project at
\url{https://github.com/quangng2000/hjepa-lean} machine-checks finite scalar
and finite-dimensional vector versions of the abstraction identities,
three-level telescoping and retained-low-level readout results,
algebraic consequences of the
stated OLS risk formulas, the conditional composed-planner bound, rollout
and temporal-grouping identities, and finite sign/control counterexamples.
It does not yet formalize the Gaussian OLS expectation formulas, the
general continuous-population probability arguments, or the linear spectral result.
The repository README provides the exact theorem-by-theorem coverage.
Analytical proofs below must not be conflated with complete Lean formalizations.

\section{Comparison and attribution}
The reference is Zhang et al.~\cite{hjepa}. Its experimental flat baseline is single-level LeWorldModel. H-JEPA learns distinct representations at progressively longer temporal strides and plans through lower-level subgoals. Experimental upper state encoders are pointwise; action encoders aggregate intervening actions. Prediction uses SIGReg, with inverse dynamics additionally used for DROID.

Appendix L already analyzes background-only slow-feature collapse, limitations of marginal regularization, conditional prediction variance, and inverse-dynamics barriers. Those are prior results, not proposed contributions here. We develop a common-target risk comparison and a conditional guarantee for the composed planner.

\paragraph{Fair comparison.}
Hold fixed the data distribution, observations, primitive-action interface, goal, physical horizon, and true evaluation cost. Distinguish training compute from planning compute, including every refinement level. A flat JEPA here means an action-conditioned world model, rather than image-only masked-feature pretraining.

Raw training losses from different learned spaces are insufficient: rescaling an unconstrained target and predictor by $\alpha$ rescales squared error by $\alpha^2$. Regularization can constrain this freedom, but matching scale does not ensure that two spaces retain the same task information.

\section{Exact abstraction gain--loss identity}
Freeze a low-level representation $Z$ and an upper representation $U=g(Z)$. Let $Y\in L^2(\R^m)$ be the \emph{same} evaluation target, such as a future physical state or task outcome. Define
\[
 \mu_Z=\E[Y\mid Z],\quad \mu_U=\E[Y\mid U],\quad
 R^*(V)=\inf_f\E\norm{Y-f(V)}^2 ,
\]
where the infimum is over all measurable predictors with finite risk.

\begin{theorem}[Information penalty]
\label{thm:information}
\[
 B:=R^*(U)-R^*(Z)=\E\norm{\mu_Z-\mu_U}^2\geq0 .
\]
Equality holds exactly when $\mu_Z$ has a version measurable with respect to $U$.
\end{theorem}
\begin{proof}
Conditional expectation minimizes squared loss. Expand
$Y-\mu_U=(Y-\mu_Z)+(\mu_Z-\mu_U)$.
The cross term vanishes since $\mu_Z-\mu_U$ is $Z$-measurable and
$\E[Y-\mu_Z\mid Z]=0$. Take expected squared norms. The nonnegative difference is zero exactly when $\mu_Z=\mu_U$ almost surely.
\end{proof}

Equality means conditional-mean sufficiency for this target and loss, not full sufficiency of the target distribution. For example, let $Z$ be a fair draw from $\{1,2\}$, $\epsilon$ an independent fair sign, $Y=Z\epsilon$, and $U=0$. Both conditional means vanish, but $|Y|$ reveals $Z$.

For a deterministic chain $U_0=Z,U_1,\ldots,U_L$, penalties telescope:
\[
 R^*(U_L)-R^*(Z)=
 \sum_{\ell=1}^{L}\E\norm{\E[Y\mid U_{\ell-1}]-\E[Y\mid U_\ell]}^2.
\]

Let $\widehat f_V$ use a training sample independent of the evaluation pair. Feature maps are fixed in advance, or learned on a separate sample. In the latter case, condition every conditional mean, Bayes risk, and risk identity in this section on the realized feature-learning sample and maps; one may subsequently average those conditional identities. Define
\[
 A_V=\E\norm{\widehat f_V(V)-\E[Y\mid V]}^2,
\]
including training randomness. Assume both learned predictors have finite expected test risk, so all differences below are well-defined. Conditional orthogonality gives
$R(\widehat f_V)=R^*(V)+A_V$, hence
\begin{equation}
 \boxed{R(\widehat f_U)-R(\widehat f_Z)=B+A_U-A_Z.}
 \label{eq:central}
\end{equation}
Abstraction helps expected prediction risk exactly when
\[
 \boxed{A_Z-A_U>B.}
\]
Learning error $A_V$ includes approximation, estimation, and optimization effects. An extra layer does not automatically reduce it.

\paragraph{Strict-loss example.}
For $Z\sim\mathcal N(0,1)$, $U=Z^2$, and $Y=Z$, symmetry gives $\mu_U=0$. Thus $R^*(Z)=0$ and $R^*(U)=1$: abstraction loses the sign.

\paragraph{Scope.}
The theorem compares nested frozen features, including levels within one trained hierarchy. A separately trained flat encoder need not equal H-JEPA's first-level encoder. If an upper level sees more history, define $Z$ to contain that same history. For action-conditioned prediction at a common physical horizon, apply the result to inputs $(Z,A)$ and $(g(Z),A)$, or $(g(Z),h(A))$ when action pooling is also part of the abstraction. A controller using both $Z$ and $U$ retains $Z$; this is not a theorem that the complete architecture loses all upper-level discarded information.

\section{Three-level H-WavJEPA specialization}
For fixed nested representations $Z_2=g(Z_1)$ and $Z_3=h(Z_2)$, define
$B_{12}=\E\norm{\mu_1-\mu_2}^2$ and
$B_{23}=\E\norm{\mu_2-\mu_3}^2$, with $\mu_\ell=\E[Y\mid Z_\ell]$.
The repository now machine-checks, on finite probability spaces with the
required finite feature-label types and scalar or finite-dimensional vector
targets, the identities
\[
 R^*(Z_3)-R^*(Z_1)=B_{12}+B_{23},\qquad
 R(f_3)-R(f_1)=B_{12}+B_{23}+A_3-A_1.
\]
Strict improvement occurs exactly when $A_1-A_3>B_{12}+B_{23}$.
The cumulative information gap vanishes exactly when both adjacent gaps
vanish. All vector losses are sums of coordinatewise squared errors for the
same target, not losses in different learned latent spaces.

Keeping the uncompressed bundle $C=(Z_1,Z_2,Z_3)$ gives
$R^*(C)=R^*(Z_1)$: a predictor of $Z_1$ lifts using the first projection,
while every predictor of $C$ composes with the deterministic bundle map to
become a predictor of $Z_1$. Lean checks these risk-preserving transformations.
This is an unrestricted-predictor statement, not a guarantee for a particular
trained readout or a subsequent lossy bottleneck.

A checked finite witness uses two independent fair signs with
$Y=S_1+2S_2$. Retaining both signs, retaining only $S_1$, and retaining
neither gives optimal risks $0,4,5$, respectively. A chosen constant-zero
high-level predictor nevertheless has lower risk ($5$) than a poor
constant-three low-level predictor ($14$). This establishes possible gain
and strict loss in a toy model, not an audio-training result.

The low/medium/high interpretation motivates an audio research question;
temporal specialization is not proved. All levels must have the same
accessible context: when an upper level pools a block, the compared lower
feature must retain the entire block. Independently trained WavJEPA and
H-WavJEPA encoders, generalization, and training convergence are outside these
theorems. The action-conditioned planning results below do not by themselves
apply to passive audio.

\section{A finite-sample example with an exact gain/loss threshold}
\begin{theorem}[Gaussian-design regression]
\label{thm:ols}
Let $Z=(U,W)\sim\mathcal N(0,I_p)$, $p=d+q$, with $U\in\R^d$, $W\in\R^q$, and
\[
 Y=U^\top a+W^\top b+\epsilon,\qquad
 \E\epsilon=0,\qquad \Var(\epsilon)=\sigma^2 .
\]
Assume independent noise $\epsilon$, iid training pairs, a fresh independent test pair, and $n>p+1$. Fit ordinary least squares \emph{without an intercept}, using either all $p$ coordinates or only $U$. Set $B=\norm b^2$. Then
\[
 R_{\rm full}=\sigma^2\frac{n-1}{n-p-1},\qquad
 R_{\rm reduced}=(\sigma^2+B)\frac{n-1}{n-d-1},
\]
and
\begin{equation}
 \boxed{R_{\rm reduced}<R_{\rm full}
 \iff B<\frac{\sigma^2q}{n-p-1}.}
 \label{eq:threshold}
\end{equation}
Equality and the reverse inequality follow by replacing $<$ accordingly.
\end{theorem}
\begin{proof}
For an $n\times r$ standard Gaussian design $X$, independent noise variance $\tau^2$, and correctly specified linear regression,
\[
 \widehat\beta-\beta=(X^\top X)^{-1}X^\top\epsilon,\qquad
 R=\tau^2+\tau^2\E\tr(X^\top X)^{-1}.
\]
The risk expression follows from independent whitened test features. The Schur complement gives
\[
 [(X^\top X)^{-1}]_{jj}=
 \{x_j^\top(I-P_{-j})x_j\}^{-1},
\]
where $P_{-j}$ projects onto the other $r-1$ columns. Conditional on those columns, the denominator is $\chi^2_{n-r+1}$. Since
$\E[(\chi^2_\nu)^{-1}]=1/(\nu-2)$ for $\nu>2$, the risk equals
$\tau^2[1+r/(n-r-1)]$.

For the full model use $r=p,\tau^2=\sigma^2$. For the reduced model, $W^\top b+\epsilon$ is independent of $U$, with variance $B+\sigma^2$, so use $r=d,\tau^2=B+\sigma^2$. Subtract:
\[
 R_{\rm reduced}-R_{\rm full}=
 \frac{(n-1)\{B(n-p-1)-\sigma^2q\}}
 {(n-d-1)(n-p-1)}.
\]
All denominator terms are positive, proving the threshold.
\end{proof}

For $n=30,p=10,d=2,q=8,\sigma^2=1$, the threshold is $8/19\approx0.421$:
\begin{center}
\begin{tabular}{rrrr}
\toprule
Discarded signal $B$ & Full risk & Reduced risk & Better predictor\\
\midrule
$0$ & $1.5263$ & $1.0741$ & Reduced\\
$0.1$ & $1.5263$ & $1.1815$ & Reduced\\
$1$ & $1.5263$ & $2.1481$ & Full\\
\bottomrule
\end{tabular}
\end{center}
These are exact toy-model expectations, not measured H-JEPA scores. With dimensions fixed and $n\to\infty$, the threshold tends to zero: discarding fixed nonzero signal eventually hurts this comparison. Removing nuisance coordinates can help finite-sample estimation, but an added nonlinear encoder does not automatically reduce effective complexity. This example does not claim that the paper shrinks latent width at every level.

\section{What a simplified predictive objective selects}
\begin{proposition}[Temporal correlation selection]
\label{prop:ar}
Let a stationary Gaussian process satisfy
\[
 X_{t+1}=D X_t+\varepsilon_{t+1},\quad
 D=\operatorname{diag}(\rho_1,\ldots,\rho_p),\quad
 1\geq\rho_1\geq\cdots\geq\rho_p\geq0 .
\]
Assume $X_t\sim\mathcal N(0,I_p)$, and temporally iid Gaussian innovations of covariance $I_p-D^2$, independent of the past. For $1\leq m\leq p$, an encoder $W$ with $WW^\top=I_m$, and integer $k\geq1$, predict $WX_{t+k}$ by $MWX_t$. Then
\[
 \min_M\E\norm{WX_{t+k}-MWX_t}^2
 =m-\norm{WD^kW^\top}_F^2,
\]
and the minimum over $W$ is $\sum_{i=1}^{m}(1-\rho_i^{2k})$, attained by selecting the first $m$ coordinates.
\end{proposition}
\begin{proof}
Encoded marginals have covariance $I_m$ and cross-covariance $C=WD^kW^\top$. Completing the square gives loss
$m-\norm C_F^2+\norm{M-C}_F^2$.
Let $P=W^\top W$. Since $I-P$ is an orthogonal projector,
\[
 \tr(WD^{2k}W^\top)-\tr\{(WD^kW^\top)^2\}
 =\norm{(I-P)D^kW^\top}_F^2\geq0.
\]
Also $\tr(WD^{2k}W^\top)\leq\sum_{i=1}^m\rho_i^{2k}$: the diagonal entries of $P$ lie in $[0,1]$ and sum to $m$, so the largest $m$ diagonal entries of $D^{2k}$ maximize the weighted sum. Coordinate selection attains both bounds.
\end{proof}

This is a tied, population-whitened linear encoder with an autonomous one-context-step predictor. It does not model neural optimization or finite-sample SIGReg. It exhibits a useful filter when selected coordinates are task-relevant and a harmful filter otherwise. For example, with $\rho_1=.9,\rho_2=.5,m=1$ and task $Y=X_{t,2}$, an optimal predictive encoder selects coordinate 1, giving task Bayes MSE $1$; retaining coordinate 2 would give $0$.

For any fixed $m$, changing stride $k$ preserves the ranking in this model. If all $\rho_i<1$, absolute explained-variance gaps tend to zero at very long horizons, although relative ratios can grow. A coordinate with $\rho_i=1$ is a random constant through time; its population variance comes from different initializations, not within-trajectory variation. Such a background can dominate the predictive objective. Fast periodic motion can remain predictable outside this AR model. Slowness, predictability, and control relevance are distinct properties.

\section{A conditional guarantee for the composed planner}
Fix initial information, a physical horizon, and a task. Let $\mathcal A$ be feasible primitive-action sequences and $J$ their common true cost, with $J^*=\inf_{\mathcal A}J$. Let $\mathcal B$ be macro-plan candidates and $\mathcal R:\mathcal B\to\mathcal A$ the actual refinement procedure. Define
\[
 \mathcal A_H=\mathcal R(\mathcal B),\qquad
 \rho_H=\inf_{a\in\mathcal A_H}J(a)-J^*\geq0.
\]
Assume nonempty sets and finite relevant infima.

\begin{theorem}[Planning regret]
\label{thm:planning}
Suppose a hierarchical surrogate $Q_H$ satisfies
\[
 \sup_{b\in\mathcal B}|Q_H(b)-J(\mathcal R(b))|\leq\delta_H
\]
and the returned $\widehat b$ satisfies
$Q_H(\widehat b)\leq\inf_bQ_H(b)+\eta_H$. Then
\begin{equation}
 \boxed{J(\mathcal R(\widehat b))-J^*\leq
 2\delta_H+\eta_H+\rho_H.}
 \label{eq:regret}
\end{equation}
A flat planner searching all of $\mathcal A$ under analogous assumptions satisfies
$J(\widehat a_F)-J^*\leq2\delta_F+\eta_F$.
\end{theorem}
\begin{proof}
\begin{align*}
 J(\mathcal R(\widehat b))
 &\leq Q_H(\widehat b)+\delta_H\\
 &\leq\inf_bQ_H(b)+\eta_H+\delta_H\\
 &\leq\inf_bJ(\mathcal R(b))+\eta_H+2\delta_H\\
 &=J^*+\rho_H+\eta_H+2\delta_H.
\end{align*}
The flat result sets $\mathcal R$ to the identity.
\end{proof}

Here $\delta_H$ includes task-cost misalignment, predictive error, and refinement error; $\eta_H$ measures optimization error of the whole composition; $\rho_H$ measures lost access to good primitive plans. A shorter search might reduce $\eta_H$, but requires proof or measurement. Local solver guarantees do not establish global $\eta_H$. Upper-level prediction loss alone does not bound $\delta_H$, and latent subgoals need not be realizable.

A smaller upper bound establishes a better guarantee, not smaller realized cost. Actual superiority needs a direct comparison or a suitable lower bound for the flat algorithm. If both surrogates are exact in true task-cost units ($\delta_H=\delta_F=0$), refinement loses no optimal value ($\rho_H=0$), and exact optimizers exist and are found ($\eta_H=\eta_F=0$), both methods attain $J^*$. Perfect latent prediction alone is insufficient.

\subsection{From rollout error to task error}
Let $T=kN$, let $u_j$ be the true abstract state after $j$ refined macro-steps, and let $\widehat u_j$ be its prediction. Assume, uniformly over relevant states and candidate plans,
\[
 \norm{u_{j+1}-\widehat G(u_j,b_j)}\leq\epsilon_H,\qquad
 \norm{\widehat G(u,b)-\widehat G(v,b)}
 \leq L_H\norm{u-v}.
\]
The first assumption includes abstract-state insufficiency, macro-action representation, dynamics error, and refinement. Triangle inequality and induction yield
\[
 e_{j+1}\leq\epsilon_H+L_He_j,\qquad
 e_N\leq L_H^Ne_0+\epsilon_HS_N(L_H),\qquad
 S_N(L)=\sum_{i=0}^{N-1}L^i .
\]
For $Q_H(b)=C_H(\widehat u_N(b))$, suppose $C_H$ is $K_H$-Lipschitz and
$|J(\mathcal R(b))-C_H(u_N(b))|\leq\kappa_H$ uniformly. Adding and subtracting the true abstract-state cost gives
\[
 \delta_H\leq\kappa_H+
 K_H\{L_H^Ne_0+\epsilon_HS_N(L_H)\}.
\]
Squared latent distance needs a bounded relevant domain for uniform Lipschitzness. Cumulative task cost requires a sum of stagewise error terms.

\begin{proposition}[Grouping steps alone gives no tighter bound]
If a macro predictor is just $k$ compositions of the flat predictor, with inherited bounds
$L_H=L_F^k$ and $\epsilon_H=\epsilon_FS_k(L_F)$, then
\[
 \epsilon_H S_N(L_H)=\epsilon_FS_{kN}(L_F).
\]
\end{proposition}
\begin{proof}
$S_k(L)S_N(L^k)=\sum_{j=0}^{N-1}\sum_{i=0}^{k-1}L^{kj+i}=S_{kN}(L)$.
\end{proof}
A better learned predictor, abstraction, or optimization problem must supply the improvement. Merely shortening the macro horizon is insufficient. These statements concern a fixed-horizon plan; receding-horizon execution needs a separate closed-loop analysis.

\section{Counterexample for control and compute accounting}
Let an observed state be $x=(c,m,y)$, with $c\sim\mathcal N(0,1)$, an independent fair sign $m$, and initially $y=0$. Choose $a\in\{-1,+1\}$, with
\[
 c'=c,\qquad m'=m,\qquad y'=\mathbf1\{a=m\}.
\]
The goal is $y'=1$. The abstraction $\phi(x)=c$ is nonconstant and standard Gaussian, and admits exact prediction $\widehat G(c,a)=c$. Nevertheless, a controller observing only $c$ and a goal providing no information about $m$ succeeds with probability $1/2$. Conditional on its observations and independent policy randomness, $m$ remains fair. A full-state controller chooses $a=m$ and succeeds surely.

Perfect nonconstant latent prediction therefore does not imply control sufficiency. The full H-JEPA may retain $m$ at a lower level; this example does not prove unavoidable architectural failure. It identifies an assumption the composed-controller guarantee must address.

\paragraph{Compute.}
For a sequential rollout solver, flat planning costs approximately $I_FTC_F$, with $I_F$ iterations and per-step forward/backward cost $C_F$. A two-level full-plan refinement scheme could cost
\[
 C_{\rm hier}\approx I_H(T/k)C_H+
 (T/k)I_{\rm low}kC_{\rm low}+C_{\rm encoding}.
\]
The refinement term still scales with $T$. Actual solver schedules determine total cost. There is no universal exponential speedup from a hierarchy alone. Added encoders and predictors may raise training cost; longer clips can reduce the number of usable training samples.

\section{Next research target}
The elementary results above establish conditional mechanisms and counterexamples, not novelty. A stronger contribution would connect $\kappa_H,\epsilon_H,\eta_H,\rho_H$ to properties enforced or estimable by H-JEPA training, rather than assuming favorable values.

\begin{enumerate}
 \item Specify a controlled stochastic environment family, true task cost, and exact flat and hierarchical algorithms.
 \item Prove when upper representations preserve task value and when macro-actions admit feasible refinements; make approximate Markov and action-sufficiency assumptions explicit.
 \item Derive finite-sample guarantees for learned representations and dynamics, accounting for temporal dependence and independent evaluation.
 \item Bound the actual multilevel optimizer or solve comparison instances exactly; guarantees for individual local objectives are insufficient.
 \item Exhibit proved gain and loss regimes under the same metric and stated compute budget, then check novelty against abstraction and hierarchical-control theory.
\end{enumerate}

Experiments should separate flat planning, multiple horizons in a shared space, upper-level costs with flat rollout, and the complete hierarchy. Match data and report training and planner compute separately. Sweep task-relevant discarded signal, stride, depth, action ambiguity, temporal dependence, and sample size. Measure real task outcomes as well as latent loss. A finite-capacity probe measures recoverability under that probe; it does not certify Bayes information loss.

Read Theorem~\ref{thm:information} and identity~\eqref{eq:central} first, then the explicit threshold~\eqref{eq:threshold}, and finally the composed-planner guarantee~\eqref{eq:regret}. The research question is when the reduction in learning or optimization difficulty exceeds information and refinement costs.

\begin{thebibliography}{9}
\bibitem{hjepa}
W. Zhang, B. Terver, M. Rabbat, Y. LeCun, and R. Balestriero.
\emph{H-JEPA: End-to-End Learning of Hierarchical World Models for Visual Planning}.
arXiv:2610.06805v1, October 5, 2026.
\url{https://arxiv.org/abs/2610.06805v1}.
Architecture and planning: \S2; existing collapse theory: Appendix L.
\end{thebibliography}
\end{document}
